\section{Consequences and open questions}\label{sec:synthesis}
The central distinction is whether locally admissible representations can
be made simultaneous. For envelope comparisons this is a probability law
with all the prescribed means. Positive coefficients permit common upper
attainment and reduce the comparison to lower-envelope rounding; signs
instead lead to cut ranges. Structural restrictions improve the available
couplings, while positive physical lower bounds change the coefficients
and may change the relevant incidence structure. In spatial certification,
a pseudoexpectation can satisfy a node's prescribed moments and polynomial
constraints without supplying a feasible graph distribution. Counting
regions that preserve it requires both a degree calculation and a witness
cover argument.

Several bounded refinements sharpen this picture. The fixed-ambient
balanced orientation improves the positive-box upper bound by replacing
two with $\beta_N<2$, and the coefficient-regularity cardinality argument
uses the same law under its stated coefficient ratios. The general-radix
cutoff construction gives exact attainment beyond the simpler large-radix
regime. Coordinatewise graph lifts admit endpoint interpolation without
polynomial-degree loss, including everywhere-finite nonpolynomial
coordinate functions. Bounded signed-monomial lifts admit the lower-bound
transfer at order $r\ge1$ when $rD\ge2$, and the quadratic formulation has
a separate order-one finite upper certificate. The unlifted cubic still
requires $r\ge2$. These completions retain their domain, oracle and
full-graph agreement hypotheses.

The remaining questions are quantitative or concern stronger methods.
The exact cubic constant lies between the certified lower value
$1610000/743033$ and $31/12$; optimality of the three-law mixture does not
determine it. The optimized harmonic upper asymptotic has no matching
second-order lower term. For a fixed physical aspect ratio the interval
$\max\{2,\rho_B\}\le C_{\rm box}(\rho_B)\le\rho_B+2$ is not closed by
the large-aspect leading constant. Exact structural constants for incidence
treewidth at least three remain open; the width-two coloring proof does
not extend merely by increasing the number of colors. On unequal-aspect
positive boxes even bipartite frequency-two exactness fails, and the
bipartite supremum remains between $3/2$ and two.

For spatial methods, arbitrary affine branching, additional coupled-cut
localizers, and convex hulls of feasible graphs exceed the proved oracle
scope. The balanced-halfspace calculation identifies an obstruction to the
preserved-moment proof under affine branching; it provides neither a fast
affine tree nor a lower bound for all such trees. The examples therefore
suggest precise questions about stronger local information, rather than
a conclusion about every global optimization algorithm.
